\chapter{Guiding Principles for Practitioners}
\label{ch:guiding-principles}

\emph{What rules does a practitioner follow before, during and after practice, and why?}
The syllabus asks for the principles, not a fixed list. Say them in your own words and give the reason for each.
An instructor follows them first and then teaches them, because a class that ignores them is where injuries and complaints start.

\section{Guiding principles for practitioners}

\subsection*{Before practice}
\begin{itemize}
\item Keep the place, the body and the mind clean.
\item Practise in a calm, quiet place.
\item Practise on an empty or light stomach. A little honey in warm water is allowed if weak.
\item Empty the bladder and bowels first.
\item Use a mat or a folded blanket.
\item Wear light, comfortable cotton clothes.
\end{itemize}

\subsection*{When not to practise}
\begin{itemize}
\item Do not practise when exhausted, ill, in a hurry or under acute stress.
\item A person with a long-term illness, pain or a heart problem should ask a physician or yoga therapist first.
\item A person who is pregnant or menstruating should ask a yoga expert first.
\end{itemize}

\begin{caution}
Do not start a student with a heart problem, a long-term illness or pain until a physician or yoga therapist has been asked. A certified teacher must review this chapter before print.
\end{caution}

\subsection*{During practice}
\begin{itemize}
\item Begin with a prayer or invocation.
\item Move slowly and in a relaxed way, aware of the body and the breath.
\item Do not hold the breath unless the practice says so.
\item Breathe through the nose unless told otherwise.
\item Never hold the body tight or jerk it.
\item Work within your own capacity.
\item Practise regularly and with patience, because results take time.
\item Keep the cautions of each practice in mind.
\item End with meditation, silence or a peace chant.
\end{itemize}

\subsection*{After practice}
Wait about 20 to 30 minutes before a bath or a meal.

\subsection*{Food}
A light, mostly vegetarian diet is the usual advice. Two meals a day can be enough for an adult who is not ill or doing heavy work. Treat this as advice, not as a rule to impose on students.

\subsection*{Five-point guide to a yogic life}
The five points are credited to Swami Vishnudevananda, developed from the teaching of Swami Sivananda (as stated by the Sivananda Yoga Vedanta Centres).
\begin{itemize}
\item Proper exercise (asana).
\item Proper breathing (pranayama).
\item Proper relaxation (shavasana).
\item Proper diet (vegetarian).
\item Positive thinking and meditation (Vedanta and dhyana).
\end{itemize}

\begin{recall}
\begin{itemize}
\item Before: clean place, body and mind; empty or light stomach; empty bladder and bowels; mat; light cotton clothes.
\item Do not practise when exhausted, ill, in a hurry or under acute stress.
\item During: prayer first; slow and aware; no breath holding unless the practice says so; within your capacity; end with meditation, silence or a peace chant.
\item After: wait about 20 to 30 minutes before a bath or a meal.
\item Diet advice is advice, not a rule to impose.
\end{itemize}
\end{recall}

\subsection*{Common questions}
\begin{itemize}
\item \emph{Must I learn the principles word for word?} No. Say them in your own words and give the reason for each.
\item \emph{Should a teacher ask before touching a student to correct a posture?} Yes, ask first. This is sound teaching practice, but the Common Yoga Protocol and the guidebook do not state it, so keep it as advice and not as a Board rule.
\item \emph{Can I take something before practice if I feel weak?} A little honey in warm water is allowed. The safe teaching line is that a full stomach makes practice uncomfortable.
\item \emph{Some slides say to wait 40 to 60 minutes and avoid water after practice. Which is right?} Use about 20 to 30 minutes before a bath or a meal, as in the Common Yoga Protocol. No source supports avoiding water.
\item \emph{Has a certified teacher reviewed this?} Not yet. It is checked against the named sources, but a certified teacher or qualified reviewer must still review it.
\end{itemize}

\subsection*{Worked examples}
\begin{itemize}
\item \textbf{Model answer, principles before, during and after practice.} Sort the points under three heads and give a reason for each. Before: clean place, body and mind; calm, quiet place; empty or light stomach; empty bladder and bowels; mat or folded blanket; light cotton clothes. During: prayer or invocation; slow and aware; no breath holding unless the practice says so; never jerk or hold the body tight; within your capacity; regular and patient; end with meditation, silence or a peace chant. After: wait about 20 to 30 minutes before a bath or a meal.
\item \textbf{A student arrives straight after lunch.} Explain that practice should be on an empty or light stomach, because a full stomach makes practice uncomfortable. Suggest moving the class to later or doing only gentle breathing. Honey in warm water is allowed only if a person is weak and the stomach is otherwise empty.
\item \textbf{A student with a heart problem, or with high blood pressure, wants to join or to learn forceful breathing.} Do not start the practice. A person with a long-term illness, pain or a heart problem should ask a physician or yoga therapist first. Point to the cautions of the practice and offer a gentler option once the physician has agreed.
\end{itemize}

\subsection*{Check yourself}
\begin{enumerate}
\item How long should one wait after practice before a bath or a meal?
\item Fill in the blank: a session begins with a prayer or invocation and ends with meditation, silence or a \underline{\hspace{2cm}}.
\item Name the five points of a yogic life credited to Swami Vishnudevananda.
\item True or false: it is fine to hold the breath in any practice if you feel strong.
\item Give two situations in which one should not practise.
\item Why does an instructor follow these principles before teaching them?
\item A pregnant student asks to join your class. What do you say?
\item A student says everyone must eat only vegetarian food and two meals a day. How do you respond?
\item May a student hold the breath during any practice?
\item What should a teacher do before touching a student to correct a posture?
\item Why practise on an empty or light stomach?
\end{enumerate}

\subsection*{Answers}
\begin{enumerate}
\item About 20 to 30 minutes.
\item Peace chant.
\item Proper exercise (asana), proper breathing (pranayama), proper relaxation (shavasana), proper diet (vegetarian), and positive thinking with meditation (Vedanta and dhyana).
\item False. Do not hold the breath unless the practice itself says so.
\item Any two of: when exhausted, ill, in a hurry or under acute stress.
\item The principles protect the practitioner and make the practice work. A class that ignores them is where injuries and complaints start.
\item A person who is pregnant or menstruating should ask a yoga expert first.
\item A light, mostly vegetarian diet is the usual advice, and two meals a day can be enough for an adult who is not ill or doing heavy work. Treat it as advice, not a rule to impose.
\item Only where the practice itself says so.
\item Ask permission first (advice, not a Board rule).
\item So that the body can use its energy for the practice and avoid discomfort.
\end{enumerate}

\paragraph{Sources.} The Board syllabus for Level 1 (Yoga Protocol Instructor), the Common Yoga Protocol general guidelines, and the official guidebook for Level I (2016, older scheme, used as a coverage check). The five-point guide is as stated by the Sivananda Yoga Vedanta Centres.
